\documentclass[12pt]{article}

\usepackage{amsmath,amssymb,amsthm}
\usepackage{mathtools}
\usepackage{tikz-cd}
\usepackage{tikz}
\usepackage{hyperref}
\usepackage{cleveref}
\usepackage{graphicx}
\usepackage[margin=1in]{geometry}
\usepackage{everypage}
\usepackage{xcolor}
\usepackage{booktabs}
\usepackage{enumitem}
\usepackage{microtype}
\usepackage{array}
\usepackage{longtable}
\usepackage{setspace}

\setstretch{1.08}
\hypersetup{
  colorlinks=true,
  linkcolor=blue!55!black,
  citecolor=blue!55!black,
  urlcolor=blue!65!black
}

\newtheorem{theorem}{Theorem}[section]
\newtheorem{proposition}[theorem]{Proposition}
\newtheorem{lemma}[theorem]{Lemma}
\newtheorem{corollary}[theorem]{Corollary}
\theoremstyle{definition}
\newtheorem{definition}[theorem]{Definition}
\newtheorem{example}[theorem]{Example}
\newtheorem{counterexample}[theorem]{Counterexample}
\theoremstyle{remark}
\newtheorem{remark}[theorem]{Remark}

\newcommand{\C}{\mathbb C}
\newcommand{\R}{\mathbb R}
\newcommand{\Z}{\mathbb Z}
\newcommand{\N}{\mathbb N}
\newcommand{\A}{\mathcal A}
\newcommand{\B}{\mathcal B}
\newcommand{\Kalg}{K_{\mathrm{alg}}}
\newcommand{\Ktop}{K^{\mathrm{top}}}
\newcommand{\ProFin}{\mathrm{ProFin}}
\newcommand{\Cont}{\operatorname{Cont}}
\newcommand{\State}{\operatorname{State}}
\newcommand{\Spec}{\operatorname{Spec}}
\newcommand{\Eq}{\operatorname{Eq}}
\newcommand{\Status}[1]{\textsc{#1}}

\definecolor{grokgray}{RGB}{110,110,110}
\AddEverypageHook{%
  \ifnum\value{page}=1
    \begin{tikzpicture}[remember picture,overlay]
      \node[
        rotate=90,
        anchor=south,
        font=\Large\sffamily\bfseries\color{grokgray},
        inner sep=0pt
      ] at ([xshift=38pt,yshift=0.52\paperheight]current page.south west)
      {GrokRxiv:2026.07.condensed-observables\quad
       [\,math.OA\,]\quad
       14 Jul 2026};
    \end{tikzpicture}
  \fi
}

\title{Condensed Families of Quantum Observables:\\
Positivity, C*-Norms, and Solid K-Theory}
\author{Matthew Long \\
\textit{The YonedaAI Collaboration} \\
\textit{YonedaAI Research Collective} \\
Chicago, IL \\
\texttt{matthew@yonedaai.com} $\cdot$ \url{https://yonedaai.com}}
\date{14 July 2026}

\begin{document}

\maketitle

\begin{abstract}
We isolate the operator algebraic layer needed to place quantum lattice systems in a
condensed mathematical framework.  For a compact Hausdorff parameter space $S$ and
a C*-algebra $A$, the section algebra $C(S,A)$ has pointwise multiplication and
involution, fiberwise positivity, and the supremum C*-norm.  We prove descent for
finite jointly surjective covers and prove that, for profinite $S$, functions through
finite clopen quotients are norm dense.  These statements justify a precise
objectwise C*-enhancement of the represented condensed algebra $\underline A$, but
they do not manufacture a norm, an order, or a state space from a bare condensed
ring.  We then connect this construction to quasi-local observable algebras,
compatible local states, and crossed products for covariant disorder.  We give a
careful account of Ko Aoki's solidification theorem: connective solidified
algebraic K-theory recovers connective operator K-theory for Banach algebras, while
periodic operator K-theory requires inversion of the complex degree-two or real
degree-eight Bott element.  Counterexamples separate continuous state families
$C(S,\State(A))$ from states of $C(S,A)$, continuous profinite families from locally
constant ones, and algebraic data from a chosen C*-completion.  The resulting bridge
is rigorous at the level of observable algebras and K-theory.  It is not an
equivalence between microscopic Hamiltonians and effective field theories.
\end{abstract}

\tableofcontents

\section{Introduction}

The operator algebra of observables is where several parts of the proposed
condensed phase program first meet.  A lattice interaction is local data.  Its
thermodynamic dynamics acts on a norm-completed quasi-local algebra.  Ground states
are positive normalized functionals on that algebra.  Free-fermion phase invariants
are K-theory classes of projections, unitaries, or symmetry-compatible flattened
Hamiltonians.  Disorder is naturally encoded by a compact configuration space and
a crossed product.  Condensed mathematics can organize continuous and profinite
families of all these objects, but only after the analytic structures have been
specified.

The distinction matters.  A ring does not remember a norm.  A *-algebra may admit
several inequivalent C*-completions.  A positive cone is not determined by ordinary
ring operations.  A state on an algebra of sections is not the same as a continuous
family of states on the fibers.  These are not technical objections around the
edges of the theory.  They determine which objects have a physical interpretation.

This paper supplies the observable-algebra component of the sequence

\[
\begin{gathered}
\text{local interactions}
\longrightarrow
\text{condensed Hamiltonian families}
\longrightarrow
\text{uniformly gapped families}
\\
\longrightarrow
\text{stable phases}
\longrightarrow
\text{invertible phase spectra}.
\end{gathered}
\]

Our results are deliberately divided by status.

\begin{itemize}[leftmargin=2em]
\item \Status{Established}: represented condensed sets, compactly parameterized
  C*-algebras, C*-descent, quasi-local inductive limits, crossed-product observable
  algebras, functional calculus, and Aoki's solidification theorem under its stated
  hypotheses.
\item \Status{Proposed}: the use of these constructions as the observable component
  of a condensed moduli stack of Hamiltonians.
\item \Status{Open}: a comparison theorem identifying the proposed condensed stack
  with a microscopic classification or with a stack of effective field theories.
\item \Status{Obstructed}: any formulation in which condensed descent alone selects
  a C*-norm, positivity, or a physical state space.
\end{itemize}

The main mathematical contributions are as follows.

\begin{enumerate}[leftmargin=2em]
\item We formulate and prove an objectwise C*-descent theorem for $S\mapsto C(S,A)$.
\item We prove the finite-clopen approximation theorem for profinite probes and make
  clear that it is a norm-density result, not a consequence of the sheaf axiom alone.
\item We distinguish external parameter families of fiber states from states on a
  section algebra by an explicit counterexample.
\item We show that compactly parameterized sections commute with the quantum-spin
  quasi-local inductive limit.
\item We state the disordered crossed-product construction and the gapped
  functional-calculus map to K-theory with exact hypotheses.
\item We state Aoki's connective and Bott-periodic comparisons separately in the
  complex and real cases.
\end{enumerate}

Clausen and Scholze provide the condensed foundations
\cite{ClausenScholze}.  Barwick and Haine provide a closely related homotopy-coherent
formalism under the name pyknotic objects \cite{BarwickHaine}.  The quasi-local
framework follows Nachtergaele, Sims, and Young \cite{NSY}.  The disordered
crossed-product viewpoint follows Bellissard, van Elst, and Schulz-Baldes
\cite{Bellissard} and its K-theoretic development in
\cite{BourneKellendonkRennie,Kellendonk}.  The exact solid K-theory comparison is due
to Aoki \cite{Aoki}.

\section{Condensed parameter objects}

\subsection{The site and represented objects}

Fix an uncountable strong limit cardinal $\kappa$ as in
\cite{ClausenScholze}.  We suppress it from notation when no confusion can result.

\begin{definition}[Condensed set]
A condensed set is a sheaf of sets on the category of $\kappa$-small profinite sets,
with finite jointly surjective families as covers.
\end{definition}

Concretely, a presheaf $X:\ProFin^{\mathrm{op}}\to\mathrm{Set}$ is a sheaf if it
sends the empty set to a point, sends finite disjoint unions to products, and
satisfies the equalizer condition

\[
X(S)\longrightarrow
\Eq\bigl(X(S')\rightrightarrows X(S'\times_S S')\bigr)
\]

for every surjection $S'\to S$ of profinite sets.

A finite jointly surjective family $\{S_i\to S\}$ gives one surjection
$\coprod_iS_i\to S$.  Since the disjoint union is compact and $S$ is Hausdorff, this
map is a quotient map.  Thus the single-cover equalizer above is exactly the Cech
equalizer for a finite covering family.

\begin{proposition}[Represented topological spaces]
For a topological space $T$, the assignment
\[
\underline T(S)=\Cont(S,T)
\]
is a condensed set.  After the cardinal convention is fixed, the functor
$T\mapsto\underline T$ is fully faithful on $\kappa$-compactly generated spaces.
\end{proposition}

\begin{proof}
Finite disjoint unions give products of mapping sets.  If $q:S'\to S$ is a
surjection, then $q$ is a quotient map because its domain is compact and its
codomain is Hausdorff.  A continuous map $S'\to T$ that agrees on
$S'\times_S S'$ therefore descends uniquely to a continuous map $S\to T$.
Full faithfulness on compactly generated spaces is Proposition 1.7 of
\cite{ClausenScholze}, with the stated cardinal convention.
\end{proof}

Every Banach space is metrizable, hence first countable and compactly generated.
Banach algebras and C*-algebras therefore lie safely in the fully faithful range.
This observation lets us use represented condensed objects without replacing the
norm topology by an unspecified topology.

\begin{definition}[Condensed anima]
A condensed anima is an appropriate hypercomplete anima-valued sheaf on compact
Hausdorff spaces, equivalently on a suitable profinite or extremally disconnected
basis.  It is the homotopy-coherent replacement for a condensed set.
\end{definition}

Condensed anima can retain automorphisms and higher equivalences.  This is relevant
for moduli of Hamiltonians, where gauge equivalences and phase automorphisms should
not be collapsed to equality.  Not every condensed anima is represented by a
topological space.  Representability is a property, not part of the definition.

\subsection{What the sheaf condition says}

The sheaf condition is a gluing statement along covers.  It does not state that a
functor preserves cofiltered limits, nor that an invariant is determined by all
finite quotients of a profinite set.  Finite-resolution approximation below will
come from the supremum norm and clopen partitions.

This distinction prevents a common overstatement.  A profinite set is an inverse
limit of finite discrete sets, but an arbitrary sheaf need not convert that inverse
limit into a direct limit.  For $C(S,A)$ there is a dense direct limit, followed by
completion.  The completion is indispensable.

\section{Compactly parameterized C*-algebras}

\subsection{Definitions and conventions}

\begin{definition}[C*-algebra]
A complex C*-algebra is a complex Banach *-algebra $A$ whose norm satisfies
\[
\|a^*a\|=\|a\|^2
\]
for every $a\in A$.
\end{definition}

For $a=a^*\in A$, the following conditions are equivalent: $a=b^*b$ for some
$b\in A$, the spectrum of $a$ lies in $[0,\infty)$, and $a$ has a positive square
root.  We write $A_+$ for the resulting closed convex cone.

\begin{definition}[Section algebra]
Let $S$ be compact Hausdorff and $A$ a C*-algebra.  Define
\[
C(S,A)=\{f:S\to A\mid f\text{ is norm-continuous}\}.
\]
Operations and involution are pointwise, and
\[
\|f\|_\infty=\sup_{s\in S}\|f(s)\|_A.
\]
\end{definition}

\begin{theorem}[Objectwise C*-structure]\label{thm:objectwise}
Let $S$ be compact Hausdorff and $A$ a complex C*-algebra.  Then $C(S,A)$ is a
C*-algebra.  It is unital if and only if $A$ is unital.  Its positive cone is
fiberwise:
\[
C(S,A)_+
=
\{f\in C(S,A):f(s)\in A_+\text{ for all }s\in S\}.
\]
There is a canonical C*-isomorphism
\[
C(S,A)\cong C(S)\otimes_{\min}A.
\]
\end{theorem}

\begin{proof}
Compactness makes the supremum norm finite.  Pointwise estimates give
$\|fg\|_\infty\leq\|f\|_\infty\|g\|_\infty$ and
$\|f^*\|_\infty=\|f\|_\infty$.  A uniform limit of continuous $A$-valued
functions is continuous, so the section algebra is complete.  Finally,
\[
\|f^*f\|_\infty
=\sup_s\|f(s)^*f(s)\|
=\sup_s\|f(s)\|^2
=\|f\|_\infty^2.
\]
This proves the C*-identity.

If $f\geq0$, each evaluation map is a *-homomorphism, hence $f(s)\geq0$.  Conversely,
suppose every $f(s)$ is positive.  The square-root map on the positive cone is norm
continuous by continuous functional calculus.  Thus $g(s)=f(s)^{1/2}$ is continuous
and $f=g^*g$.  The tensor-product identification is the standard map on elementary
tensors, $h\otimes a\mapsto(s\mapsto h(s)a)$.  Since $C(S)$ is commutative and
nuclear, its minimal and maximal tensor products agree.  The minimal completion
therefore gives all continuous $A$-valued functions without a representation-dependent
choice of tensor norm.
\end{proof}

\begin{corollary}[Functoriality]
A continuous map $u:T\to S$ induces a contractive *-homomorphism
\[
u^*:C(S,A)\longrightarrow C(T,A),\qquad f\longmapsto f\circ u.
\]
It is unital when $A$ is unital, and it is isometric when $u$ is surjective.
\end{corollary}

\begin{proof}
All algebraic statements are pointwise.  The norm inequality follows because the
supremum is taken over a subset of the values of $f$.  Surjectivity makes the two
sets of values equal.
\end{proof}

\subsection{C*-descent}

\begin{theorem}[Objectwise C*-descent]\label{thm:descent}
Let $\{q_i:S_i\to S\}_{i=1}^n$ be a finite jointly surjective family of maps between
profinite sets.  For every C*-algebra $A$, restriction induces an isometric
C*-isomorphism
\[
C(S,A)
\cong
\Eq\left(
\prod_i C(S_i,A)
\rightrightarrows
\prod_{i,j}C(S_i\times_S S_j,A)
\right),
\]
where finite products carry the maximum norm.
\end{theorem}

\begin{proof}
Let $q:\coprod_iS_i\to S$ be the jointly surjective map.  It is a quotient map.  A
compatible tuple $(f_i)$ defines a function on the disjoint union that is constant
on the fibers of $q$, so it descends uniquely to a continuous function $f:S\to A$.
This proves the equalizer statement at the level of sets.  Pointwise operations show
that the map and its inverse are *-homomorphisms.  Joint surjectivity gives
\[
\|f\|_\infty=\max_i\|f_i\|_\infty,
\]
so the isomorphism is isometric.
\end{proof}

\begin{corollary}[Locality of order and norm]
Under the hypotheses of \cref{thm:descent}, a section is positive if and only if all
of its pullbacks are positive.  Its norm is the maximum of the pullback norms.
\end{corollary}

This is the exact statement in which positivity and the C*-norm are compatible with
condensed descent.  The norm and order are inherited from $A$ before descent is
applied.  Descent does not construct them from a ring.

\subsection{A condensed observable-algebra object}

For a C*-algebra $A$ with its norm topology, define
\[
\underline A(S)=C(S,A).
\]
By \cref{thm:objectwise,thm:descent}, this is a sheaf whose values are C*-algebras
and whose restriction maps are contractive *-homomorphisms.  We call it the
\emph{represented objectwise C*-algebra sheaf}.  This term avoids presupposing a
universal formal category named condensed C*-algebras.

For the phase program, this construction is sufficient.  It provides compact and
profinite parameter families, a norm topology for gapped homotopies, fiberwise
positivity, and functorial pullback.  More elaborate categorical definitions can be
added later, but the present theorems do not depend on them.

\section{Finite clopen approximation}

\begin{definition}[Finite clopen quotient]
For a profinite set $S$, a finite clopen partition $P$ determines a finite discrete
quotient $q_P:S\to P$.  A function $S\to A$ factors through $P$ exactly when it is
constant on every part of the partition.
\end{definition}

The finite clopen partitions form a directed set under refinement.  If $A$ is a
C*-algebra, then $C(P,A)\cong A^P$ with the maximum norm.

\begin{theorem}[Finite-resolution density]\label{thm:clopen}
Let $S$ be profinite and $A$ a Banach space.  Then
\[
C(S,A)
=
\overline{
\bigcup_{P} C(P,A)
}^{\|\cdot\|_\infty},
\]
where the union ranges over finite clopen quotients of $S$.  When $A$ is a
C*-algebra, this is a dense *-subalgebra and its completion is $C(S,A)$.
\end{theorem}

\begin{proof}
Fix $f\in C(S,A)$ and $\varepsilon>0$.  For each $s\in S$, continuity gives a
neighborhood on which the oscillation of $f$ is smaller than $\varepsilon$.  Since
$S$ is zero dimensional, the neighborhood can be chosen clopen.  Compactness gives
a finite clopen cover, and Boolean refinement gives a finite clopen partition $P$
with oscillation smaller than $\varepsilon$ on every part.  Choose one point $s_U$
in each $U\in P$ and define $f_P|_U=f(s_U)$.  Then
$\|f-f_P\|_\infty<\varepsilon$.  Pointwise operations preserve locally constant
functions, proving the C*-algebra statement.
\end{proof}

\begin{counterexample}[Continuous does not mean locally constant]\label{ce:cantor}
Let $S=\{0,1\}^{\N}$ and define
\[
f(x)=\sum_{n=1}^{\infty}2^{-n}x_n.
\]
The series converges uniformly, so $f:S\to\R$ is continuous.  It does not factor
through any finite quotient because changing a sufficiently late coordinate changes
its value.  Its prefix truncations are locally constant and converge uniformly.
\end{counterexample}

The counterexample identifies the role of completion.  Finite disorder data provide
approximants.  The completed C*-algebra contains limits that depend on infinitely
many coordinates.

\begin{remark}[Sheaf descent versus inverse limits]
The sheaf condition used in \cref{thm:descent} concerns compatible data over a cover.
\Cref{thm:clopen} concerns approximation by a directed family of finite quotients.
Neither theorem implies the other.  Their combination is useful because it provides
both gluing and controlled finite-resolution approximation.
\end{remark}

\section{Norm and positivity are not algebraic}

\subsection{The missing structure problem}

Suppose one begins with a condensed ring $R$.  Ring operations do not specify a
topology, a complete norm, an involution, or a positive cone.  Adding an involution
still does not select a C*-norm.  The next example is elementary and decisive.

\begin{counterexample}[Multiple C*-norms]\label{ce:norms}
Give $\C[x]$ the involution determined by $x^*=x$.  For each compact infinite
$K\subset\R$, define
\[
\|p\|_K=\sup_{t\in K}|p(t)|.
\]
This is a faithful pre-C*-norm.  The norms obtained from $K=[-1,1]$ and
$K=[-2,2]$ differ, and their completions are $C([-1,1])$ and $C([-2,2])$.
Their induced orders differ as well: $1-x^2$ is positive in the first completion but
not in the second.
\end{counterexample}

\begin{proof}
A nonzero polynomial cannot vanish on an infinite compact subset of $\R$, so the
seminorm is faithful.  Pointwise complex conjugation gives
\[
\|p^*p\|_K=\sup_{t\in K}|p(t)|^2=\|p\|_K^2.
\]
The Stone-Weierstrass theorem identifies the completions.  Finally,
$1-t^2\geq0$ on $[-1,1]$ but takes negative values on $[-2,2]$.
\end{proof}

The example remains an obstruction after applying the represented condensed functor.
The same algebraic *-object can inherit distinct topological and ordered
enhancements.  Therefore the correct input to the observable construction is a
chosen C*-algebra, not a ring from which one hopes to recover a physical completion.

\subsection{Faithful representations and physical choice}

A faithful *-representation $\pi:A\to B(\mathcal H)$ of a pre-C*-algebra supplies
the operator norm $\|a\|=\|\pi(a)\|$.  Different representation classes can lead to
different universal or reduced completions.  Crossed products provide a familiar
example when the acting group is not amenable.  In the disorder models below the
acting group is $\Z^d$, which is amenable, so full and reduced crossed products agree.

The physical choice still includes the algebra of observables and its representation.
A norm alone does not specify a ground state, a GNS Hamiltonian, or a thermodynamic
spectral gap.

\section{States and parameter families}

\begin{definition}[State]
For a unital C*-algebra $A$, a state is a complex linear map
$\omega:A\to\C$ such that
\[
\omega(1)=1,
\qquad
\omega(a^*a)\geq0
\]
for every $a\in A$.
\end{definition}

The state space $\State(A)$ carries the weak-* topology.  It is compact and convex.
For a compact parameter space $S$, there are two natural but different constructions.

\begin{enumerate}[leftmargin=2em]
\item $C(S,\State(A))$ is the space of weak-* continuous families
  $s\mapsto\omega_s$ of states on a fixed fiber algebra.
\item $\State(C(S,A))$ is the state space of the C*-algebra of sections.
\end{enumerate}

\begin{counterexample}[The two state constructions differ]\label{ce:states}
Let $A=\C$ and let $S$ contain at least two points.  Then $\State(\C)$ is a point, so
$C(S,\State(\C))$ is a point.  By the Riesz representation theorem,
$\State(C(S))$ is the compact convex space of regular Borel probability measures on
$S$.  It contains every point mass and their nontrivial convex combinations.
\end{counterexample}

The difference has a physical interpretation.  A continuous family $\omega_s$
assigns a state at each externally fixed parameter.  A state on $C(S,A)$ may also
average over the parameter.  Under separability and standard Borel hypotheses, its
description involves a probability measure and measurable fiber functionals.  That
extra measure is not part of a continuous family of fiber states.

\begin{proposition}[Evaluation states]
Let $A$ be unital, $s\in S$, and $\omega\in\State(A)$.  Then
\[
\operatorname{ev}_{s,\omega}(f)=\omega(f(s))
\]
is a state on $C(S,A)$.
\end{proposition}

\begin{proof}
Evaluation at $s$ is a unital *-homomorphism.  Composing a state with a unital
*-homomorphism gives a state.
\end{proof}

Evaluation states form only a subclass.  Convex combinations and integrals of them
already produce many additional states.  Consequently, a moduli problem must say
whether $S$ is an external parameter, a disorder variable to be averaged, or part of
the observable algebra.

\section{Quasi-local observable algebras}

\subsection{The local net}

Let $\Gamma$ be a countable set, usually a uniformly locally finite metric lattice.
At each site $x\in\Gamma$, fix a finite-dimensional Hilbert space $\mathcal H_x$.
For a finite region $\Lambda\Subset\Gamma$, define
\[
\A_\Lambda
=
\bigotimes_{x\in\Lambda}B(\mathcal H_x).
\]
Finite dimensionality removes tensor-norm ambiguity.  If
$\Lambda\subset\Lambda'$, define
\[
\iota_{\Lambda,\Lambda'}(a)
=
a\otimes1_{\Lambda'\setminus\Lambda}.
\]

\begin{theorem}[Quasi-local inductive limit]\label{thm:quasilocal}
Each inclusion $\iota_{\Lambda,\Lambda'}$ is unital, multiplicative, and isometric.  The
algebraic union
\[
\A_{\Gamma}^{\mathrm{loc}}
=
\bigcup_{\Lambda\Subset\Gamma}\A_\Lambda
\]
has a well-defined C*-norm, and its completion
\[
\A_\Gamma
=
\overline{\A_{\Gamma}^{\mathrm{loc}}}^{\|\cdot\|}
=
\varinjlim_{\Lambda\Subset\Gamma}\A_\Lambda
\]
is the quasi-local observable algebra.
\end{theorem}

\begin{proof}
An injective *-homomorphism between C*-algebras is isometric.  The maps are
compatible, so the norm on an element of the union is independent of the region in
which it is represented.  Completing the directed union gives the C*-inductive
limit.  This is the standard construction used in \cite{NSY}.
\end{proof}

\begin{proposition}[Local approximation of positivity]\label{prop:positive-local}
The positive cone of the quasi-local algebra satisfies
\[
(\A_\Gamma)_+
=
\overline{
\bigcup_{\Lambda\Subset\Gamma}(\A_\Lambda)_+
}.
\]
\end{proposition}

\begin{proof}
The right side lies in $(\A_\Gamma)_+$ because *-homomorphisms preserve positivity
and the positive cone is closed.  Conversely, write $a=b^*b$ with
$b\in\A_\Gamma$.  Choose local $b_n\to b$.  Then $b_n^*b_n$ are local positive
elements and converge to $b^*b$.
\end{proof}

\begin{proposition}[Compatible local states]\label{prop:local-states}
Suppose states $\omega_\Lambda\in\State(\A_\Lambda)$ satisfy
\[
\omega_{\Lambda'}\circ\iota_{\Lambda,\Lambda'}=\omega_\Lambda
\]
whenever $\Lambda\subset\Lambda'$.  Then they extend uniquely to a state
$\omega\in\State(\A_\Gamma)$.  Every state on $\A_\Gamma$ restricts to such a
compatible family.
\end{proposition}

\begin{proof}
Compatibility defines a linear functional of norm one on the algebraic union.
Positivity holds because each element of the local union belongs to some finite
algebra.  The functional extends uniquely by continuity to the norm completion.
The converse follows by restriction.
\end{proof}

\subsection{Compact parameter spaces commute with completion}

\begin{theorem}[Parameterized quasi-local limit]\label{thm:param-limit}
Let $S$ be compact Hausdorff.  Then
\[
C(S,\A_\Gamma)
\cong
\overline{
\bigcup_{\Lambda\Subset\Gamma}C(S,\A_\Lambda)
}^{\|\cdot\|_\infty}.
\]
The isomorphism preserves involution, positivity, and the C*-norm.
\end{theorem}

\begin{proof}
By \cref{thm:objectwise},
\[
C(S,\A_\Gamma)\cong C(S)\otimes_{\min}\A_\Gamma.
\]
The commutative algebra $C(S)$ is nuclear.  Minimal tensoring by it commutes with
injective C*-inductive limits.  Applying this to \cref{thm:quasilocal} proves the
density statement.  The structural claims follow from the isometric
*-isomorphism.
\end{proof}

For profinite $S$, one can see the theorem without tensor-product machinery.  A
continuous section has compact image.  Approximate finitely many values by elements
of a common local algebra, then use a sufficiently fine clopen partition to form a
locally constant local-algebra-valued approximation.

\subsection{Interactions are not global bounded Hamiltonians}

An interaction is a map
\[
\Phi:\{X\Subset\Gamma\}\longrightarrow\A_X^{\mathrm{sa}}.
\]
The finite-volume Hamiltonian is
\[
H_\Lambda=\sum_{X\subseteq\Lambda}\Phi(X).
\]
Even when the interaction has rapid decay, the sum over the infinite lattice need
not converge in $\A_\Gamma$.  Under the locality hypotheses of \cite{NSY}, the
thermodynamic object is a strongly continuous automorphism group
$\tau_t:\A_\Gamma\to\A_\Gamma$ or its generator on a dense domain.

This observation is important for any Hamiltonian moduli stack.  Infinite-volume
interactions, finite-volume Hamiltonians, dynamics, ground-state representations,
and bounded observables are related but distinct objects.

\section{Disorder and crossed products}

\subsection{Profinite configuration spaces}

Let $D$ be a finite discrete set and define
\[
\Omega=D^{\Z^d}.
\]
With the product topology, $\Omega$ is compact, totally disconnected, and metrizable,
hence profinite.  The translation action
\[
\alpha:\Z^d\curvearrowright\Omega
\]
is continuous.  Cylinder functions, which depend on finitely many coordinates, are
locally constant and dense in $C(\Omega)$ by \cref{thm:clopen}.

\subsection{The covariant bulk algebra}

\begin{definition}[Crossed-product observable algebra]
For a compact space $\Omega$ with a continuous $\Z^d$ action $\alpha$, define
\[
\B=C(\Omega)\rtimes_\alpha\Z^d.
\]
In the presence of magnetic translations with cocycle $\theta$, define the twisted
crossed product
\[
\B_\theta=C(\Omega)\rtimes_{\alpha,\theta}\Z^d.
\]
\end{definition}

The algebra is generated by coefficient functions and translation unitaries with
the covariance relation
\[
U_a f U_a^*=\alpha_a(f).
\]
Since $\Z^d$ is amenable, its full and reduced crossed products agree.  This removes
a completion ambiguity that would be present for a general nonamenable group.

\begin{proposition}[Finite-range covariant families]
A bounded, finite-range, covariant one-particle Hamiltonian with $N$ internal orbitals
defines a self-adjoint element of
\[
M_N\bigl(C(\Omega)\rtimes_{\alpha,\theta}\Z^d\bigr).
\]
Conversely, the finite-support algebraic crossed product describes such finite-range
covariant operators, and its C*-completion includes norm limits of them.
\end{proposition}

\begin{proof}[Proof sketch]
Write the hopping operator as a finite sum
\[
h=\sum_{a\in F}h_aU_a,
\]
where $F\subset\Z^d$ is finite and $h_a\in M_N(C(\Omega))$.  Covariance is encoded by
the crossed-product relation, while self-adjointness gives the corresponding relation
between $h_a$ and $h_{-a}$.  Such finite sums form the algebraic crossed product,
which is dense by definition in the C*-crossed product.  Magnetic phases modify the
multiplication by $\theta$.  This is the framework used in
\cite{Bellissard,BourneKellendonkRennie}.
\end{proof}

The crossed product directly models covariant one-particle systems.  It does not
replace the quasi-local algebra for a general interacting spin system.  An
interacting disordered model needs a disorder action on interactions, a quasi-local
dynamics, and a specified state or representation.

\section{Gapped functional calculus and K-theory}

\subsection{From a Hamiltonian to a projection}

\begin{theorem}[Uniformly gapped functional calculus]\label{thm:gapped-k}
Let $B$ be a unital C*-algebra.  Let
\[
t\longmapsto h_t\in M_n(B)_{\mathrm{sa}}
\]
be norm-continuous for $t\in[0,1]$.  Suppose there is $\delta>0$ such that
\[
\Spec(h_t)\cap(-\delta,\delta)=\varnothing
\]
for every $t$.  Then
\[
p_t=\chi_{(-\infty,0)}(h_t)
\]
is a norm-continuous path of projections in $M_n(B)$.  Consequently,
$[p_t]\in K_0(B)$ is independent of $t$.
\end{theorem}

\begin{proof}
Choose a continuous function $g:\R\to[0,1]$ that equals $1$ on
$(-\infty,-\delta]$ and $0$ on $[\delta,\infty)$.  Since the spectrum avoids the
interpolation interval, continuous functional calculus gives
$g(h_t)=\chi_{(-\infty,0)}(h_t)$.  Functional calculus is continuous for
norm-continuous self-adjoint families, so $p_t$ is norm-continuous.  Homotopic
projections represent the same $K_0$ class.
\end{proof}

The theorem gives a controlled map
\[
h\longmapsto p_h\longmapsto[p_h]\in K_0(B).
\]
For a chiral system, spectral flattening gives an off-diagonal unitary and an odd
K-class.  Anti-linear symmetries require real, graded, or equivariant variants.
Kellendonk develops the role of graded real structures in this classification
\cite{Kellendonk}.  Bourne, Kellendonk, and Rennie relate the resulting classes to
crossed-product K-homology and bulk-edge pairings
\cite{BourneKellendonkRennie}.

\subsection{What a K-class forgets}

The class $[p_h]$ is stable under matrix stabilization and norm-continuous gapped
homotopy.  It does not retain the eigenvalue magnitudes removed by flattening.  It
does not retain a preferred interaction decomposition, Lieb-Robinson velocity,
correlation length, state, boundary condition, or response coefficient.  Response
coefficients arise only after pairing a K-class with a trace, cyclic cocycle,
K-homology class, or Kasparov class.

Equality of K-classes can be weaker than the microscopic equivalence relation.  A
classification theorem must specify the algebra, symmetry class, stabilization, and
allowed homotopies.  Outside a controlled free-fermion sector, no universal
identification is asserted here.

\section{Solidification and operator K-theory}

\subsection{The condensed algebraic K-theory object}

Let $A$ be a real or complex Banach algebra.  Regard it as a represented condensed
algebra by
\[
\underline A(S)=C(S,A).
\]
Applying algebraic K-theory produces a condensed spectrum whose value is
schematically
\[
S\longmapsto\Kalg(C(S,A)).
\]
Solidification is a completion operation on condensed spectra introduced in the
Clausen-Scholze program.  Aoki proves that it recovers established topological
K-theories in this setting \cite{Aoki}.

\begin{theorem}[Aoki, complex connective comparison]\label{thm:aoki-complex}
Let $A$ be a complex Banach algebra.  The solidification of its connective algebraic
K-theory is discrete and canonically equivalent to the connective part of operator
K-theory:
\[
\bigl(K^{\mathrm{cn}}(\underline A)\bigr)^{\mathrm{sol}}
\simeq
\bigl(\tau_{\geq0}\Ktop(A)\bigr)^{\mathrm{disc}}.
\]
\end{theorem}

\begin{proof}[Source and proof architecture]
This is Theorem A of \cite{Aoki}.  Aoki first proves discreteness under openness of
the general linear groups, a condition satisfied by Banach algebras.  The connective
comparison is characterized by agreement on $K_0$, homotopy invariance under
$A\mapsto C([0,1],A)$, and excision for the specified Banach-algebra pullback
squares.  These properties identify the solidified connective theory with
$\tau_{\geq0}\Ktop$.
\end{proof}

\begin{theorem}[Aoki, complex Bott inversion]\label{thm:aoki-bott}
Let $A$ be a complex Banach algebra and let
\[
\beta_{\C}\in K^{\mathrm{top}}_2(\C)\cong\Z
\]
be a Bott generator.  Then
\[
\Kalg(\underline A)^{\mathrm{sol}}[\beta_{\C}^{-1}]
\simeq
\Ktop(A)^{\mathrm{disc}}.
\]
\end{theorem}

\begin{proof}
This is the periodic part of Theorem A in \cite{Aoki}.  The connective comparison
provides the nonnegative homotopy groups.  Inverting the degree-two complex Bott
class imposes complex Bott periodicity and yields the periodic operator K-theory
spectrum.
\end{proof}

\begin{theorem}[Aoki, real comparison]\label{thm:aoki-real}
For a real Banach algebra $A$, solidified connective algebraic K-theory identifies
with connective real operator K-theory.  If
$\beta_{\R}$ is the real Bott class in degree $8$, then
\[
\Kalg(\underline A)^{\mathrm{sol}}[\beta_{\R}^{-1}]
\simeq
\Ktop(A)^{\mathrm{disc}}.
\]
\end{theorem}

\begin{proof}[Source and qualification]
Aoki proves the real version in the comparison section of \cite{Aoki}.  The real
period is eight, not two.  The theorem uses the real Banach-algebra category and the
real Bott element.  Complexification does not justify silently replacing this class
by the complex degree-two generator.
\end{proof}

The localization at $\beta_{\R}$ is defined through the module-spectrum structure
over $\Kalg(\underline{\R})^{\mathrm{sol}}$.  Bott inversion is therefore a
localization in spectra, not an informal operation on homotopy groups.

\subsection{The semitopological comparison}

Aoki also proves that if $A$ is a complex associative algebra equipped with the
condensed structure induced by the ordinary topology on $\C$, then
solidification of its algebraic K-theory recovers semitopological K-theory.  This is
separate from the Banach-algebra operator K-theory statement.  The hypotheses and
target should not be interchanged.

\subsection{Exact nonclaims}

The notation
\[
\Kalg(\underline A)^{\mathrm{sol}}
\simeq
\Ktop(A)
\]
without further qualification is generally too strong.  The exact statement needs
either connective truncation or Bott inversion.  Even after the exact equivalence is
applied, it compares K-theory spectra.  It does not reconstruct

\begin{itemize}[leftmargin=2em]
\item the C*-norm on $A$,
\item the positive cone $A_+$,
\item the state space $\State(A)$,
\item the observable algebra from its K-theory,
\item a local interaction that generates the dynamics,
\item or an effective field theory.
\end{itemize}

This is still a substantial bridge.  It shows that a central invariant of operator
algebraic phase theory arises through an intrinsic condensed completion of algebraic
K-theory.  The bridge is invariant-level, exact, and narrower than an equivalence of
physical models.

\section{A proposed observable component of the phase stack}

\subsection{Input data}

We now state the new organizational proposal.  It is not needed for the proofs above.
Fix a lattice $\Gamma$, finite-dimensional on-site Hilbert spaces, a symmetry type,
and a class of interactions with specified locality estimates.  For a profinite test
object $S$, an observable-family datum consists of

\begin{enumerate}[leftmargin=2em]
\item a continuous family of interactions with the required uniform locality bound,
\item the fixed quasi-local C*-algebra $\A_\Gamma$,
\item the objectwise section algebra $C(S,\A_\Gamma)$,
\item optional weak-* continuous fiber states $S\to\State(\A_\Gamma)$,
\item and symmetry actions by continuous *-automorphisms.
\end{enumerate}

If disorder changes the covariant observable algebra, the datum may instead use an
appropriate crossed product.  The choice must be part of the object.  It cannot be
inferred from the condensed parameter set.

\begin{definition}[Proposed observable functor]
Let $\mathfrak{Obs}_{\Gamma}$ assign to a profinite $S$ the groupoid of the above
observable-family data and structure-preserving equivalences.  Pullback along
$T\to S$ is defined by restriction of the parameter family.
\end{definition}

\Status{Proposed}: after the interaction and morphism categories are specified,
$\mathfrak{Obs}_{\Gamma}$ should be enhanced to a condensed anima or stack.
\Cref{thm:descent} proves the descent statement for its fixed section-algebra
component.  It does not prove descent for gaps, ground states, dynamics, or phase
equivalences.  Those require separate analytic theorems.

\subsection{Morphisms}

The morphism class determines the moduli problem.  Possible morphisms include
star isomorphisms, covariant star homomorphisms, Morita equivalences, quasi-local
automorphisms, and stable gapped homotopies.  They are not interchangeable.

For the present paper, we use contractive star homomorphisms for section-algebra
functoriality and isometric star isomorphisms for equivalence.  A later phase stack may
invert a larger class after proving that the relevant physical data are preserved.

\subsection{Dependency boundaries}

The observable layer receives locality and dynamics from the locality paper.  It
supplies the norm, order, functional calculus, state language, and K-theory target to
the gap and microscopic-to-effective papers.  It does not prove any of the following:

\begin{itemize}[leftmargin=2em]
\item a uniform thermodynamic gap,
\item stability of that gap under perturbations,
\item equivalence of a lattice system and an effective field theory,
\item or microscopic realization of an abstract bordism class.
\end{itemize}

\section{Examples and counterexamples}

\subsection{A finite parameter set}

Let $S=\{0,1\}$ and $A=M_n(\C)$.  Then
\[
C(S,A)\cong A\oplus A,
\qquad
\|(a_0,a_1)\|=\max\{\|a_0\|,\|a_1\|\}.
\]
The positive cone is $A_+\oplus A_+$.  A continuous family of states is simply a
pair $(\omega_0,\omega_1)$.  A state on $A\oplus A$ also chooses a convex weight:
\[
\varphi(a_0,a_1)
=
\lambda\omega_0(a_0)+(1-\lambda)\omega_1(a_1),
\qquad 0\leq\lambda\leq1.
\]
This finite example already displays the distinction in \cref{ce:states}.
The coefficients $\lambda$ and $1-\lambda$ are the classical probability measure on
$S$.  That measure is absent from the bare parameter family $S\to\State(A)$.

\subsection{A profinite spin family}

Let $S=\{0,1\}^{\N}$ and let $A=M_2(\C)$.  A norm-continuous family of one-site
observables is a function $f:S\to M_2(\C)$.  It can depend on infinitely many bits,
but \cref{thm:clopen} approximates it uniformly by functions of finitely many bits.
Positivity is checked matrix by matrix in every fiber.  The approximation theorem
preserves positivity after applying the continuous positive-part map or square root.

This is a useful finite-resolution model for disorder.  It is not a proof that a
thermodynamic invariant computed on finite cylinders converges.  Such convergence
requires continuity of the invariant in the chosen topology.

\subsection{A spin-chain algebra}

For $\Gamma=\Z$ and $\mathcal H_x=\C^2$, each local algebra is a matrix algebra
$M_{2^{|\Lambda|}}(\C)$.  The quasi-local algebra is the norm completion of their
directed union.  Local density matrices give compatible states when their partial
traces agree.  A translation-invariant product state is obtained from a single
one-site density matrix.

No infinite sum of identical on-site Hamiltonians belongs to the unital quasi-local
algebra in norm.  Its finite-volume sums generate compatible local dynamics, and the
thermodynamic automorphism group is obtained under the standard locality hypotheses.

\subsection{A disordered covariant model}

For binary disorder $D=\{0,1\}$ on $\Z^d$, the coefficient algebra $C(\Omega)$ is
the completion of cylinder functions.  A nearest-neighbor covariant Hamiltonian is a
finite sum in $M_N(C(\Omega))\rtimes\Z^d$.  If the Fermi energy lies in a uniform
spectral gap, \cref{thm:gapped-k} gives a projection in the matrix crossed product
and hence a K-class.

If only a mobility gap is present, the spectral projection may belong to a smoother
Sobolev-type algebra rather than the C*-algebra in the same way.  Mobility-gap
invariants require extra localization hypotheses and are outside the theorem stated
here.

\section{Computational and formal interfaces}

\subsection{Finite Haskell demonstrations}

The companion program implements exact finite contracts for the following tasks:

\begin{enumerate}[leftmargin=2em]
\item adjoint involution on $2\times2$ complex matrices,
\item the principal-minor positivity criterion for Hermitian $2\times2$ matrices,
\item state evaluation $\operatorname{Tr}(\rho a)$,
\item the supremum of Hermitian spectral norms over a finite parameter set,
\item and the error bound for finite-prefix clopen approximation on
  $\{0,1\}^{\N}$.
\end{enumerate}

The program is compiled with \texttt{-Wall -Wextra -Werror}.  It provides executable
regression checks for definitions.  Floating-point tests do not prove positivity in
an abstract C*-algebra, and finite prefix tests do not prove the sheaf or completion
theorems.

\subsection{Lean representation boundary}

A Lean library can directly represent finite-dimensional *-algebras, positive
matrices, continuous functions into normed algebras, and finite clopen partitions.
It can state the quasi-local inductive system and Aoki comparison as interfaces with
provenance.  It should not introduce axioms labeled as proofs of Aoki's theorem or of
thermodynamic dynamics.  Every imported analytic theorem must retain its hypotheses.

The most useful formal dependency chain is

\[
\begin{tikzcd}[column sep=small,cells={font=\small}]
\text{local matrix algebras}
  \arrow[r]
& \text{C*-inductive limit}
  \arrow[r]
& C(S,\A_\Gamma)
  \arrow[r]
& \Ktop(C(S,\A_\Gamma)).
\end{tikzcd}
\]

The solidification theorem supplies a comparison to the last object only after the
connective or Bott-periodic qualification has been selected.

\section{Claim ledger and limitations}

\begin{longtable}{>{\raggedright\arraybackslash}p{0.17\textwidth}
                  >{\raggedright\arraybackslash}p{0.49\textwidth}
                  >{\raggedright\arraybackslash}p{0.24\textwidth}}
\toprule
Status & Claim & Dependency \\
\midrule
\endfirsthead
\toprule
Status & Claim & Dependency \\
\midrule
\endhead
\Status{Established}
& $C(S,A)$ is a C*-algebra with fiberwise positivity and supremum norm.
& Compact Hausdorff $S$, C*-algebra $A$. \\
\Status{Established}
& Objectwise C*-descent holds for finite jointly surjective profinite covers.
& Quotient-map descent. \\
\Status{Established}
& Finite clopen quotient functions are dense for profinite $S$.
& Banach target and compactness. \\
\Status{Established}
& Quantum-spin local algebras have a quasi-local C*-inductive limit.
& Chosen local tensor products and isometric embeddings. \\
\Status{Established}
& Covariant one-particle disorder is modeled by a crossed product.
& Compact hull and continuous action. \\
\Status{Established}
& A uniformly gapped norm-continuous path gives a constant K-class.
& Continuous functional calculus. \\
\Status{Established}
& Solidification recovers connective, then Bott-periodic, operator K-theory.
& Aoki's Banach-algebra hypotheses. \\
\Status{Proposed}
& These objects form the observable component of a condensed Hamiltonian stack.
& A defined moduli category and descent for all added data. \\
\Status{Open}
& The resulting stack is equivalent to a microscopic phase or EFT stack.
& No general comparison theorem. \\
\Status{Obstructed}
& A bare condensed ring canonically determines norm, order, and states.
& \Cref{ce:norms,ce:states}. \\
\bottomrule
\end{longtable}

The main limitations are structural.  We restrict the quasi-local construction to
finite-dimensional on-site Hilbert spaces.  Infinite-dimensional site algebras and
unbounded on-site terms require domain and topology choices such as those analyzed
in \cite{NSY}.  We treat true spectral gaps for bounded one-particle Hamiltonians,
not mobility gaps.  We do not classify noninvertible topological order.  We do not
derive an effective field theory or a response action from a K-class.  We do not
claim that all algebraic K-classes have local gapped microscopic representatives.

\section{Discussion}

The observable-algebra layer reveals both the promise and the limit of the condensed
perspective.  It is promising because ordinary compact parameters, profinite
disorder, norm completions, and algebraic K-theory can be placed in a common
functorial setting.  The object $S\mapsto C(S,A)$ is not a metaphor.  It satisfies
descent as a C*-algebra-valued construction, and for profinite $S$ it has a concrete
finite-resolution approximation.

The limit is equally clear.  Condensed mathematics organizes topology already
present in $A$.  It does not decide which representation gives the physical norm,
which cone is positive, which states are admissible, or whether a dynamics has a
thermodynamic gap.  These are analytic inputs.  The framework becomes useful when it
keeps them visible instead of absorbing them into a slogan.

Aoki's theorem sharpens the bridge.  Operator K-theory can be obtained from
solidified algebraic K-theory of the represented Banach algebra.  The result explains
why condensed methods can recover a central phase invariant.  It also marks a clean
boundary.  K-theory is much less information than the C*-algebra with its order,
states, locality, and dynamics.

For the wider research program, the next step is to combine the observable functor
with a uniformly local interaction functor and a uniformly gapped subfunctor.  The
descent of $C(S,A)$ is already established here.  Descent of uniform gaps and
quasi-adiabatic equivalences is a separate problem.  Any synthesis should retain
that separation.

\section{Conclusion}

We have proved the operator algebraic statements needed for a rigorous condensed
formulation of parameterized observables.  Compact families form C*-algebras of
sections.  Their norms and positive cones are fiberwise.  They satisfy descent over
profinite covers.  Profinite families admit norm-controlled finite-clopen
approximations.  Quantum-spin observables arise as C*-inductive limits, and compact
families commute with that completion.  Covariant disorder leads to crossed
products, while a uniform spectral gap gives a stable K-class by functional
calculus.

The strongest condensed-to-operator bridge is Aoki's solidification theorem, with
connective truncation before periodicity and Bott inversion for the periodic theory.
The complex Bott class has degree two and the real Bott class degree eight.  These
qualifications are part of the theorem.

The resulting picture is concise:
\[
\begin{gathered}
\text{chosen C*-observable algebra}
\longmapsto
\text{condensed family of section algebras}
\\
\longmapsto
\text{solidified algebraic K-theory}
\longmapsto
\text{operator K-theory}.
\end{gathered}
\]
It is a proved invariant-level route.  It is not a reconstruction of physical
positivity from algebra, and it is not a universal microscopic-to-effective
equivalence.

\appendix

\section{Detailed checks for the section algebra}

This appendix records elementary estimates used implicitly in
\cref{thm:objectwise}.  They are included to make the analytic assumptions auditable.

\begin{lemma}[Continuity of pointwise operations]
If $f,g\in C(S,A)$, then $f+g$, $fg$, and $f^*$ are norm-continuous.
\end{lemma}

\begin{proof}
Addition and involution are norm-continuous operations on $A$.  For multiplication,
at a fixed $s_0$ write
\[
f(s)g(s)-f(s_0)g(s_0)
=
(f(s)-f(s_0))g(s)+f(s_0)(g(s)-g(s_0)).
\]
Continuity of $g$ gives a local bound on $\|g(s)\|$, and both terms tend to zero as
$s\to s_0$.
\end{proof}

\begin{lemma}[Completeness]
Every Cauchy sequence in $C(S,A)$ for the supremum norm converges to an element of
$C(S,A)$.
\end{lemma}

\begin{proof}
If $(f_n)$ is Cauchy in the supremum norm, then $(f_n(s))$ is Cauchy in $A$ for every
$s$.  Completeness of $A$ defines $f(s)=\lim_nf_n(s)$.  The convergence is uniform.
A uniform limit of continuous maps into a metric space is continuous, so
$f\in C(S,A)$ and $f_n\to f$ in the supremum norm.
\end{proof}

\begin{lemma}[Continuity of square roots]
If $f\in C(S,A)$ is pointwise positive, then $s\mapsto f(s)^{1/2}$ is norm-continuous.
\end{lemma}

\begin{proof}
The union of the spectra of $f(s)$ is contained in
$[0,\|f\|_\infty]$.  Polynomial approximation of the square-root function on this
compact interval is uniform.  Applying the same polynomials to $f(s)$ shows that the
square-root map is the uniform limit of continuous $A$-valued functions.
\end{proof}

\begin{lemma}[Closedness of the positive cone]
The set $C(S,A)_+$ is norm closed.
\end{lemma}

\begin{proof}
If $f_n\geq0$ and $f_n\to f$ uniformly, then $f_n(s)\to f(s)$ in $A$ for every $s$.
The positive cone $A_+$ is norm closed, so $f(s)\geq0$.  Apply
\cref{thm:objectwise}.
\end{proof}

\section{Local states and density matrices}

For a finite-dimensional local algebra $\A_\Lambda=B(\mathcal H_\Lambda)$, every
state has the form
\[
\omega_\Lambda(a)=\operatorname{Tr}(\rho_\Lambda a)
\]
for a unique density matrix $\rho_\Lambda\geq0$ with trace one.  If
$\Lambda\subset\Lambda'$, compatibility of states is equivalent to
\[
\rho_\Lambda
=
\operatorname{Tr}_{\Lambda'\setminus\Lambda}(\rho_{\Lambda'}).
\]

\begin{proposition}[Product-state compatibility]
Fix one-site density matrices $\rho_x$ for $x\in\Gamma$.  The finite-volume density
matrices
\[
\rho_\Lambda=\bigotimes_{x\in\Lambda}\rho_x
\]
define a compatible family and therefore a state on $\A_\Gamma$.
\end{proposition}

\begin{proof}
The partial trace of a tensor product over the discarded factors multiplies by the
traces of those factors.  Each trace is one, so the remaining density matrix is
$\rho_\Lambda$.  Apply \cref{prop:local-states}.
\end{proof}

This construction gives a state on the observable algebra.  It does not assert that
the state is a ground state of a specified interaction.  Ground-state conditions
involve the dynamics or local energy inequalities.

\section{Counterexample audit}

The three principal counterexamples answer distinct possible overclaims.

\begin{enumerate}[leftmargin=2em]
\item \Cref{ce:norms} shows that algebraic *-data do not choose a C*-completion or
  order.
\item \Cref{ce:states} shows that a continuous family of fiber states is not the
  state space of the section algebra.
\item \Cref{ce:cantor} shows that a continuous function on a profinite set need not
  factor through a finite quotient.
\end{enumerate}

Together they force three completions or choices to remain explicit: the operator
norm completion, the distinction between parameter and averaging, and the norm
closure of finite-resolution data.

\begin{thebibliography}{99}

\bibitem{Aoki}
K. Aoki,
\emph{(Semi)topological K-theory via solidification},
arXiv:2409.01462, 2024,
\url{https://arxiv.org/abs/2409.01462}.

\bibitem{BarwickHaine}
C. Barwick and P. Haine,
\emph{Pyknotic objects, I: Basic notions},
arXiv:1904.09966, 2019,
\url{https://arxiv.org/abs/1904.09966}.

\bibitem{Bellissard}
J. Bellissard, A. van Elst, and H. Schulz-Baldes,
\emph{The non-commutative geometry of the quantum Hall effect},
Journal of Mathematical Physics 35 (1994), 5373 to 5451,
\url{https://arxiv.org/abs/cond-mat/9411052}.

\bibitem{BourneKellendonkRennie}
C. Bourne, J. Kellendonk, and A. Rennie,
\emph{The K-theoretic bulk-edge correspondence for topological insulators},
Annales Henri Poincar\'e 18 (2017), 1833 to 1866,
\url{https://arxiv.org/abs/1604.02337}.

\bibitem{ClausenScholze}
D. Clausen and P. Scholze,
\emph{Lectures on Condensed Mathematics},
University of Bonn lecture notes, 2019,
\url{https://www.math.uni-bonn.de/people/scholze/Condensed.pdf}.

\bibitem{Kellendonk}
J. Kellendonk,
\emph{On the C*-algebraic approach to topological phases for insulators},
Journal of Physics A: Mathematical and Theoretical 50 (2017), 085202,
\url{https://arxiv.org/abs/1509.06271}.

\bibitem{NSY}
B. Nachtergaele, R. Sims, and A. Young,
\emph{Quasi-locality bounds for quantum lattice systems, Part I: Lieb-Robinson
bounds, quasi-local maps, and spectral flow automorphisms},
Journal of Mathematical Physics 60 (2019), 061101,
\url{https://arxiv.org/abs/1810.02428}.

\end{thebibliography}

\end{document}
